\documentclass[11pt]{article}
\usepackage{amsmath}
\begin{document}

\section*{Step 45: Discrete RT-Consistency Constraint}

Let \(c_e\) be the variable capacities of the finite Step-42 graph. For each
boundary region \(A_i\), linearizing the unperturbed min-cut gives
\[
\delta \mathrm{Area}(A_i)=\sum_e M_{ie}\delta c_e.
\]
RT consistency under perturbation requires
\[
M\,\delta c = \delta S.
\]
In this carrier, \(M\) has 38 rows and 14
columns, with rank 10. Hence the cokernel dimension is
28, so \(P_\perp\delta S=0\) is a nontrivial
consistency condition.

The RT-preserving control \(\delta S=M\delta c\) has residual
9.2513644228e-17. The generic contracted-state
perturbation has residual 0.0418931924111, so it
does not admit an RT-consistent capacity response on this fixed graph.

This is a recognition-level finite analogue of the linearized-Einstein
consistency condition, not a dynamical equation, not a continuum tensor
equation, and not a local bulk equation of motion.

The matrix \(M\) is built from the unperturbed min-cut edge incidence. This is
a fixed-cut-chamber linearization: it is valid only while the active minimal
surface does not change. If a perturbation changes the active min-cut pattern,
\(M\) must be recomputed.

\end{document}
